\section{Loss Function}

Durante il training, la \textit{Detection Head} produce tre tipi di
predizioni:

\[
\text{classification},\qquad
(l,t,r,b),\qquad
\text{centerness}.
\]

Il \textit{Target Generator} produce gli stessi tipi di target a partire dalle
bounding box del ground truth. La funzione di loss misura quindi quanto le
predizioni della rete differiscono dai target.

Nel progetto la loss totale è definita come combinazione di tre componenti:

\[
L_{\text{total}}
=
\lambda_{\text{cls}}L_{\text{cls}}
+
\lambda_{\text{reg}}L_{\text{reg}}
+
\lambda_{\text{ctr}}L_{\text{ctr}}.
\]

Nel caso attuale tutti i pesi sono impostati a $1$, quindi le tre componenti
contribuiscono in ugual misura alla loss totale.

\subsection{Classification Loss: Focal Loss}

La classification branch deve determinare, per ogni location della feature
map, se è presente un oggetto oppure se la location appartiene al background.
Nel nostro caso è presente una sola classe, quindi il target è binario:

\[
y =
\begin{cases}
1 & \text{location positiva}\\
0 & \text{background}.
\end{cases}
\]

Per questo motivo viene utilizzata una \textit{sigmoid focal loss}.

La Focal Loss è stata introdotta da Lin et al. per affrontare il forte
sbilanciamento tra esempi positivi e negativi tipico dei \textit{dense
detectors}.\footnote{Tsung-Yi Lin et al., \textit{Focal Loss for Dense Object
Detection}, ICCV 2017.}

La prima parte è la Binary Cross Entropy:

\[
BCE(p,y)
=
-\left[
y\log(p)
+
(1-y)\log(1-p)
\right],
\]

dove $p$ è la probabilità predetta dalla rete.

La Focal Loss introduce poi un fattore che riduce il contributo degli esempi
facili:

\[
FL(p_t)
=
-\alpha_t(1-p_t)^\gamma\log(p_t).
\]

Dove $p_t$ rappresenta la probabilità assegnata alla classe corretta:

\[
p_t=
\begin{cases}
p & \text{se } y=1\\
1-p & \text{se } y=0.
\end{cases}
\]

Il termine

\[
(1-p_t)^\gamma
\]

è il termine fondamentale della Focal Loss. Se un esempio è classificato
correttamente con alta confidenza, allora $p_t$ è vicino a $1$ e il termine
diventa molto piccolo. In questo modo gli esempi facili contribuiscono poco alla
loss, mentre gli esempi difficili ricevono un peso maggiore.

Il parametro $\alpha_t$ permette invece di bilanciare il contributo delle
classi positive e negative.

Nel progetto vengono utilizzati:

\[
\alpha=0.25,\qquad\gamma=2.
\]

La Focal Loss viene utilizzata sulla classification branch per tutte le
location della feature map.

\subsection{Regression Loss: Generalized IoU}

La regression branch ha il compito di localizzare la bounding box. Per ogni
location positiva il modello predice:

\[
(l,t,r,b),
\]

ovvero le distanze della location dai quattro lati della bounding box.

Nel progetto la predizione viene prima trasformata in una bounding box
$XYXY$ e successivamente confrontata con la bounding box ground truth tramite
la \textit{Generalized Intersection over Union} (GIoU).

La IoU tra due bounding box $A$ e $B$ è:

\[
IoU =
\frac{|A\cap B|}
{|A\cup B|}.
\]

Il valore è compreso tra $0$ e $1$ e misura quanto le due bounding box
si sovrappongono.

La sola IoU presenta però un problema quando le bounding box non si
sovrappongono: in quel caso la IoU è $0$ e non fornisce informazioni sulla
distanza tra i due box.

La Generalized IoU introduce quindi il più piccolo bounding box $C$ che
contiene entrambe le box:

\[
GIoU =
IoU -
\frac{|C\setminus(A\cup B)|}
{|C|}.
\]

Il secondo termine penalizza le situazioni in cui le due bounding box sono
lontane tra loro.

La loss utilizzata è:

\[
L_{\text{GIoU}}=1-GIoU.
\]

Di conseguenza, quando la box predetta coincide con il ground truth:

\[
GIoU\rightarrow1
\qquad\Rightarrow\qquad
L_{\text{GIoU}}\rightarrow0.
\]

Nel nostro codice la regression loss viene calcolata solo sulle location
positive.

È importante notare che il modello non confronta direttamente i quattro valori
$LTRB$ con una loss L1. La procedura è:

\[
\text{LTRB prediction}
\rightarrow
\text{decode}
\rightarrow
\text{XYXY prediction}
\rightarrow
\text{GIoU loss}.
\]

Questa scelta differisce dal paper sulla pneumonia, che utilizza una
\textit{Smooth L1 loss} per la regressione.\footnote{Linghua Wu et al.,
\textit{Pneumonia detection based on RSNA dataset and anchor-free deep learning
detector}, Scientific Reports, 2024.}

\subsection{Centerness Loss}

La \textit{centerness branch} deve imparare quanto una location sia vicina al
centro della bounding box associata.

Il Target Generator calcola il target di centerness come:

\[
c^*
=
\sqrt{
\frac{\min(l,r)}{\max(l,r)}
\cdot
\frac{\min(t,b)}{\max(t,b)}
}.
\]

Il rapporto

\[
\frac{\min(l,r)}{\max(l,r)}
\]

misura quanto la location sia centrata orizzontalmente, mentre

\[
\frac{\min(t,b)}{\max(t,b)}
\]

misura quanto sia centrata verticalmente.

La radice quadrata combina i due valori.

Una location esattamente al centro del box tende ad avere:

\[
c^*\approx1,
\]

mentre una location molto vicina a uno dei bordi tende ad avere:

\[
c^*\approx0.
\]

Questo target viene utilizzato per insegnare alla centerness branch a
distinguere le location che producono box di buona qualità da quelle lontane
dal centro.

Nel progetto la centerness prediction viene confrontata con il target tramite
la stessa \textit{sigmoid focal loss} utilizzata per la classificazione:

\[
L_{\text{ctr}}
=
-\alpha_t(1-p_t)^\gamma\log(p_t).
\]

La loss di centerness viene però calcolata solamente sulle location positive.

\subsection{Normalizzazione della Loss}

Le tre componenti vengono inizialmente sommate su tutte le immagini e tutti i
livelli della FPN.

Successivamente la loss viene normalizzata rispetto al numero totale di
location positive:

\[
N_{\text{pos}}
=
\sum_{\text{images}}
\sum_{\text{levels}}
\#\{\text{positive locations}\}.
\]

La loss finale diventa quindi:

\[
L_{\text{total}}
=
\frac{
\lambda_{\text{cls}}L_{\text{cls}}
+
\lambda_{\text{reg}}L_{\text{reg}}
+
\lambda_{\text{ctr}}L_{\text{ctr}}
}
{\max(1,N_{\text{pos}})}.
\]

La normalizzazione impedisce che la scala della loss dipenda semplicemente dal
numero di location positive presenti nel batch.

\subsection{Confronto con i paper}

La loss del progetto combina idee provenienti principalmente dalla famiglia
FCOS e dalla Focal Loss, ma non coincide esattamente con la formulazione di uno
dei due paper di riferimento.

\subsubsection{Paper sulla pneumonia}

Il paper sulla pneumonia utilizza due componenti principali:

\[
L =
L_{\text{center}}
+
L_{\text{scale}}.
\]

Il CENTER viene addestrato tramite Focal Loss, mentre la previsione della scala
viene addestrata tramite Smooth L1.\footnote{Linghua Wu et al.,
\textit{Pneumonia detection based on RSNA dataset and anchor-free deep learning
detector}, Scientific Reports, 2024.}

La formulazione è quindi:

\[
\text{CENTER}
\rightarrow
\text{Focal Loss}
\]

\[
\text{SCALE}
\rightarrow
\text{Smooth L1}.
\]

\subsubsection{FCOS}

FCOS utilizza invece una formulazione basata su tre predizioni:

\[
\text{classification},
\qquad
\text{regression},
\qquad
\text{centerness}.
\]

La classification branch utilizza Focal Loss, la regressione utilizza una
loss basata sulla IoU e la centerness utilizza una Binary Cross Entropy.
Inoltre, la center-ness viene introdotta specificamente per ridurre il
contributo delle prediction generate da location lontane dal centro degli
oggetti.\footnote{Zhi Tian et al., \textit{FCOS: Fully Convolutional One-Stage
Object Detection}, ICCV 2019.}

\subsubsection{Progetto}

Il progetto è quindi concettualmente più vicino a FCOS:

\[
\boxed{
\text{Classification}
+
\text{LTRB Regression}
+
\text{Centerness}
}
\]

ma presenta alcune modifiche:

\[
\text{Classification}
\rightarrow
\text{Focal Loss}
\]

\[
\text{Regression}
\rightarrow
\text{GIoU Loss}
\]

\[
\text{Centerness}
\rightarrow
\text{Focal Loss}.
\]

La differenza più evidente rispetto a FCOS è quindi la loss della
centerness: FCOS utilizza Binary Cross Entropy, mentre nel progetto viene
utilizzata la Focal Loss.

Rispetto al paper sulla pneumonia, invece, la differenza è maggiore, perché il
progetto mantiene una classification branch esplicita, utilizza la
regressione $LTRB$ tipica della formulazione FCOS e sostituisce la Smooth L1
con la GIoU loss.

Nel complesso, la loss implementata nel progetto può quindi essere considerata
una formulazione \textit{FCOS-like}, adattata alle esigenze specifiche del
progetto.


\subsubsection{Scelta di una formulazione FCOS-like}

Sebbene il progetto sia nato con l'obiettivo di riprodurre il detector descritto
nel paper sulla pneumonia, durante l'implementazione alcune difficoltà di
ottimizzazione hanno motivato l'adozione di una formulazione più vicina a FCOS.

In particolare, nelle prime versioni del modello il training risultava
instabile, con episodi di gradient explosion. Anche una formulazione più vicina
a FCOS mostrava inizialmente problemi di stabilità. L'analisi delle
architetture e delle tecniche utilizzate nei detector ha portato all'impiego
della \textit{Group Normalization} nelle detection heads. Questa scelta è
coerente con il paper FCOS, che utilizza Group Normalization nelle convolutional
layers delle heads proprio per rendere il training più stabile.\footnote{Zhi
Tian et al., \textit{FCOS: Fully Convolutional One-Stage Object Detection},
ICCV 2019.}

La formulazione finale del progetto è quindi maggiormente vicina a FCOS rispetto
al detector descritto nel paper sulla pneumonia. Vengono infatti mantenuti
diversi elementi caratteristici di FCOS: la predizione per location, la
regressione $LTRB$, la classification branch, la center-ness branch, i
regression ranges associati ai livelli della FPN e la struttura $P_3$--$P_7$.

Il progetto non costituisce tuttavia una riproduzione esatta di FCOS. La
detection head è più semplice, poiché utilizza una sola convoluzione $3\times3$
per ciascuna tower invece delle quattro convoluzioni previste dalla
formulazione originale. Inoltre, le detection heads non condividono i parametri
tra i diversi livelli della FPN.

Anche le funzioni di loss presentano alcune differenze. La regressione utilizza
la Generalized IoU loss, coerente con una delle successive migliorie introdotte
nel lavoro FCOS, mentre la center-ness viene addestrata tramite Focal Loss
anziché Binary Cross Entropy.

Il modello può quindi essere considerato una variante semplificata e adattata
al problema di pneumonia detection di un detector \textit{FCOS-like}.


\subsubsection{Feature-map Locations and Reference Boxes}

In FCOS, ogni posizione spaziale della feature map rappresenta una
\textit{location} sulla quale il modello può effettuare una predizione. La
location viene riportata nelle coordinate dell'immagine originale utilizzando
lo stride del livello della FPN:

\[
x_{\text{img}}=(x+0.5)s,
\qquad
y_{\text{img}}=(y+0.5)s,
\]

dove $(x,y)$ è la posizione sulla feature map e $s$ è lo stride del livello.

Ad esempio, una location di $P_3$ con stride $8$ corrisponde ad un punto
dell'immagine originale distante $8$ pixel dalla location adiacente. In
termini intuitivi, una singola location può quindi essere associata ad una
cella di circa $8\times8$ pixel nell'immagine originale, anche se la feature
in quella posizione può contenere informazioni provenienti da una regione
molto più ampia a causa del \textit{receptive field} della rete.

Nel FCOS originale la location viene concettualmente trattata come un punto
rispetto al quale vengono predette direttamente le quattro distanze
\[
(l,t,r,b).
\]

Nel nostro progetto, invece, per eseguire il decoding viene utilizzato il
\texttt{BoxLinearCoder} di \texttt{torchvision}. Questo componente utilizza un
\textit{reference box} anziché un semplice punto. Per questo motivo, a partire
dalla location nell'immagine viene costruito un piccolo box centrato sulla
location e con dimensione pari allo stride:

\[
B_{\text{ref}} =
\left[
x-\frac{s}{2},
y-\frac{s}{2},
x+\frac{s}{2},
y+\frac{s}{2}
\right].
\]

Il reference box non rappresenta un'ipotesi sulla dimensione dell'oggetto e
non costituisce un anchor box nel senso tradizionale dei detector
anchor-based. È solamente una rappresentazione della location utilizzata come
riferimento dal \texttt{BoxLinearCoder}.

La pipeline di decoding diventa quindi:

\[
\text{feature-map location}
\rightarrow
\text{image coordinates}
\rightarrow
\text{reference box}
\rightarrow
\text{LTRB prediction}
\rightarrow
\text{BoxLinearCoder}
\rightarrow
\text{XYXY box}.
\]